\documentclass[11pt]{article}
\usepackage[margin=1in]{geometry}
\usepackage{amsmath,amssymb}
\usepackage{hyperref}
\usepackage{enumitem}
\usepackage{xcolor}
\usepackage{booktabs}

\hypersetup{colorlinks=true,linkcolor=blue,urlcolor=blue,citecolor=blue}

\title{How CMFGEN Works --- and the Non-Monotonic Velocity Solver}
\author{A simple explainer for first-year PhD students}
\date{}

\begin{document}
\maketitle

\section*{The big picture}

CMFGEN (CMF GENeral) is a code that computes the spectrum of light coming
out of an expanding atmosphere --- a supernova, a stellar wind, or any
outflowing envelope.  You give it the structure of the gas at one
instant in time --- the radius, density, velocity, temperature and
chemical composition at each point --- and it tells you:

\begin{itemize}[nosep]
  \item what the emergent spectrum looks like, and
  \item what the ionisation and excitation state of every atom is at
        every depth.
\end{itemize}

The hard part is that the radiation field and the atomic populations
depend on \emph{each other}.  The radiation field determines which atoms
are ionised and excited; the ionisation and excitation determine the
opacity and emissivity, which in turn determine the radiation field.  So
CMFGEN has to solve a coupled problem self-consistently.  It does this by
\emph{iteration}.

\section{Radiative transfer terms: a glossary}
\label{sec:rt-terms}

The rest of this document uses a standard set of radiative-transfer
quantities.  This section defines each one, states its physical meaning,
and notes the units and sign conventions used inside CMFGEN.  Readers
already comfortable with the notation can skip to
\S\ref{sec:two-equations}.

\subsection{Specific intensity \texorpdfstring{$I_\nu$}{I\_nu}}

The \textbf{specific intensity} (or \emph{specific radiative intensity},
sometimes called \emph{brightness}) is the most fundamental quantity in
radiative transfer.  It is the energy flowing in a given direction, per
unit area, per unit time, per unit frequency, per unit solid angle:
\begin{equation}
  I_\nu(\mathbf{r},\hat{\mathbf{n}},t)
  \;=\;
  \frac{dE}{dA_\perp\,dt\,d\nu\,d\Omega},
\end{equation}
with SI units $\mathrm{erg\,s^{-1}\,cm^{-2}\,Hz^{-1}\,sr^{-1}}$.
$I_\nu$ depends on position $\mathbf{r}$, direction $\hat{\mathbf{n}}$,
frequency $\nu$ and time $t$.  In CMFGEN the time dependence is frozen
out (a snapshot), and the geometry is spherical, so the direction is
parametrised by the ray cosine $\mu$ (see below) and the impact
parameter $p$ (see below).

The intensity is what the formal solver integrates along each ray.  It
is also the quantity that the non-monotonic solver must keep finite and
non-negative: a negative or NaN $I_\nu$ propagates through every moment
and every rate equation.

\subsection{Ray cosine \texorpdfstring{$\mu$}{mu} and impact parameter \texorpdfstring{$p$}{p}}

In spherical geometry a ray is labelled by its \textbf{impact parameter}
$p$, the closest approach distance to the centre, and by the local
\textbf{ray cosine}
\begin{equation}
  \mu \;=\; \hat{\mathbf{n}}\cdot\hat{\mathbf{r}}
  \;=\; \cos\theta,
\end{equation}
where $\theta$ is the angle between the ray direction and the outward
radial direction.  Along a given ray, $\mu$ varies with radius as
$\mu = \sqrt{1 - p^{2}/r^{2}}$; $\mu = +1$ is an outward radial ray,
$\mu = -1$ is an inward radial ray, and $\mu = 0$ is a tangential ray at
$p = r$.

CMFGEN discretises the rays by a set of impact parameters
$p_1, \dots, p_{N_P}$ (typically a Gauss quadrature on $[0,R_{\max}]$)
and integrates the transfer equation along each one.  The angle
quadrature weights turn the ray intensities into the angular moments.

\subsection{Opacity \texorpdfstring{$\chi_\nu$}{chi\_nu}}

The \textbf{opacity} (or \emph{extinction coefficient}) is the
probability per unit path length that a photon of frequency $\nu$ is
removed from the beam (by absorption or scattering out of the
direction/frequency bin).  It has units of $\mathrm{cm^{-1}}$ and is
related to the cross section $\sigma_\nu$ and number density $n$ by
$\chi_\nu = n\,\sigma_\nu$ (summed over all species).  In CMFGEN
$\chi_\nu$ includes:
\begin{itemize}[nosep]
  \item bound--bound (line) opacity;
  \item bound--free (photoionisation) opacity;
  \item free--free (bremsstrahlung) opacity;
  \item electron (Thomson) scattering opacity.
\end{itemize}
The opacity is set by the atomic populations, which is why it couples
back to the rate equations.

\subsection{Emissivity \texorpdfstring{$\eta_\nu$}{eta\_nu}}

The \textbf{emissivity} is the energy added to the beam per unit volume,
time, frequency and solid angle:
$[\eta_\nu] = \mathrm{erg\,s^{-1}\,cm^{-3}\,Hz^{-1}\,sr^{-1}}$.  It
includes thermal (free--free, free--bound) emission, bound--bound
(spontaneous + stimulated) emission, and scattering into the beam.  In
the transfer equation $\eta_\nu$ is the source term and $\chi_\nu I_\nu$
is the sink term.

\subsection{Source function \texorpdfstring{$S_\nu$}{S\_nu}}

The \textbf{source function} is the ratio of emissivity to opacity,
\begin{equation}
  S_\nu \;=\; \frac{\eta_\nu}{\chi_\nu},
\end{equation}
with the same units as $I_\nu$.  In local thermodynamic equilibrium
(LTE) $S_\nu = B_\nu(T)$, the Planck function.  In non-LTE (the general
CMFGEN case) $S_\nu$ differs from $B_\nu$ because the populations are
not at their LTE values and because scattering contributes to
$\eta_\nu$ and $\chi_\nu$ without being tied to the local $T$.

In the comoving-frame formal solver CMFGEN works with a
\textbf{generalised source function}
$S' = S + Q\,I_{\rm upwind}$ (see \S\ref{sec:two-equations}), which
absorbs the velocity-coupling contribution into an effective source.

\subsection{Optical depth \texorpdfstring{$\tau_\nu$}{tau\_nu}}

The \textbf{optical depth} along a path is the dimensionless integral
of the opacity,
\begin{equation}
  \tau_\nu \;=\; \int \chi_\nu\,ds .
\end{equation}
A photon has probability $e^{-\tau_\nu}$ of traversing the path without
interaction.  CMFGEN also uses a \textbf{generalised optical depth}
$\hat\tau = \hat\chi\,ds$, where $\hat\chi$ is the generalised opacity
that includes the velocity-coupling term (see below).  The non-monotonic
solver exists precisely because $\hat\chi$ can go negative in the
standard formulation when the velocity field is non-monotonic.

\subsection{Mean intensity \texorpdfstring{$J_\nu$}{J\_nu}}

The \textbf{mean intensity} is the angle average of $I_\nu$,
\begin{equation}
  J_\nu \;=\; \frac{1}{4\pi}\oint I_\nu\,d\Omega
  \;=\; \frac{1}{2}\int_{-1}^{+1} I_\nu(\mu)\,d\mu ,
\end{equation}
with the same units as $I_\nu$.  $J_\nu$ is the quantity that enters the
radiative rates in the statistical-equilibrium equations: a
bound--bound rate is proportional to the profile-averaged $\bar J$
(see below), and a photoionisation rate is proportional to
$J_\nu$ integrated over the ionisation edge.  Computing $J_\nu$ from
the ray intensities is the main job of the formal solver, and it is the
quantity that the non-monotonic solver replaces.

\subsection{First moment \texorpdfstring{$H_\nu$}{H\_nu} (the flux)}

The \textbf{first moment} (or \emph{Eddington flux}) is the
angle-weighted average,
\begin{equation}
  H_\nu \;=\; \frac{1}{4\pi}\oint \mu\, I_\nu\,d\Omega
  \;=\; \frac{1}{2}\int_{-1}^{+1} \mu\, I_\nu(\mu)\,d\mu .
\end{equation}
$4\pi H_\nu$ is the \emph{net flux} (energy per unit area per unit time
per unit frequency) through a spherical shell.  The luminosity is
$L_\nu = 16\pi^{2} r^{2} H_\nu$.  At a velocity reversal $H_\nu$ passes
through zero (inward and outward flux cancel), which is why the
non-monotonic configuration requires \texttt{N\_ON\_J} rather than
\texttt{G\_ONLY}: the ratio $N/H$ would be singular.

\subsection{Second moment \texorpdfstring{$K_\nu$}{K\_nu}}

The \textbf{second moment} is
\begin{equation}
  K_\nu \;=\; \frac{1}{4\pi}\oint \mu^{2}\, I_\nu\,d\Omega
  \;=\; \frac{1}{2}\int_{-1}^{+1} \mu^{2}\, I_\nu(\mu)\,d\mu .
\end{equation}
$K_\nu$ measures the radiation pressure ($P_{\rm rad} = (4\pi/c)K_\nu$)
and the degree of anisotropy.  The ratio
\begin{equation}
  f_{\rm Edd} \;=\; \frac{K_\nu}{J_\nu}
\end{equation}
is the \textbf{Eddington factor}, which closes the moment hierarchy.
In the diffusion limit $f_{\rm Edd} \to 1/3$ (isotropic); in the
free-streaming limit $f_{\rm Edd} \to 1$ (radial).  CMFGEN computes
$f_{\rm Edd}$ from the ray quadrature and feeds it into the moment
equations (or, in the non-monotonic path, into the BA matrix).

\subsection{Third moment \texorpdfstring{$N_\nu$}{N\_nu}}

The \textbf{third moment} is
\begin{equation}
  N_\nu \;=\; \frac{1}{4\pi}\oint \mu^{3}\, I_\nu\,d\Omega
  \;=\; \frac{1}{2}\int_{-1}^{+1} \mu^{3}\, I_\nu(\mu)\,d\mu .
\end{equation}
It appears in the variable-Eddington-factor closure and in the
comoving-frame moment equations (the term that couples $H$ to $K$ and
$N$).  The ratio $N_\nu/J_\nu$ is used in the \texttt{N\_ON\_J} closure
that the non-monotonic solver requires.

\subsection{Profile-averaged mean intensity \texorpdfstring{$\bar J$}{Jbar}}

For a bound--bound transition $i \to j$ the radiative rate is not
proportional to $J_\nu$ at a single frequency but to the mean intensity
\emph{averaged over the line profile} $\phi(\nu)$:
\begin{equation}
  \bar J_{ij} \;=\; \int_0^\infty J_\nu\,\phi_{ij}(\nu)\,d\nu .
\end{equation}
In the Sobolev / large-velocity-gradient limit the profile is a delta
function at the resonance frequency, so $\bar J$ reduces to $J$ at the
single resonance point.  In the comoving-frame formal solution the
profile is resolved on the frequency grid, and $\bar J$ is computed by
integrating $J_\nu$ over the line.  The Sobolev routines
(\texttt{SOBEW}, \texttt{SOBJBAR}, etc.) compute $\bar J$ directly from
the Sobolev optical depth and escape probability; these are the
routines that the Dessart+2015 sigma guard stabilises for
non-monotonic flow.

\subsection{Velocity-gradient parameter \texorpdfstring{$\sigma$}{sigma}}

The \textbf{velocity-gradient parameter} is the dimensionless logarithmic
gradient of the expansion velocity,
\begin{equation}
  \sigma \;=\; \frac{d\ln V}{d\ln R} - 1
  \;=\; \frac{R}{V}\,\frac{dV}{dR} - 1 .
\end{equation}
For homologous expansion ($V \propto R$) $\sigma = 0$.  For a
decelerating shell ($dV/dR < 0$) $\sigma < -1$, which is the regime that
breaks the standard solver.  CMFGEN maintains two arrays:
\texttt{SIGMA\_TRUE} (the unclamped physical value, used by the
non-monotonic formal solver) and \texttt{SIGMA} (clamped to
$\ge -0.999$, used by every other routine for stability).

\subsection{Sobolev optical-depth scale \texorpdfstring{$\Gamma$}{Gamma}}

$\Gamma$ is the comoving-frame frequency-gradient scale that appears in
the generalised opacity,
\begin{equation}
  \Gamma \;=\; \frac{V}{R}\bigl(1 + \sigma\,\mu^{2}\bigr) ,
\end{equation}
with units of $\mathrm{s^{-1}}$.  It measures how fast a photon's
comoving frequency shifts per unit path length.  In monotonic flow
$\Gamma > 0$ (photons redshift outward); in compression zones
$\Gamma < 0$ (photons blueshift as they move into the compression).  The
sign of $\Gamma$ is what the non-monotonic solver uses to choose the
upwind frequency.

\subsection{Generalised opacity \texorpdfstring{$\hat\chi$}{hat-chi}}

In the comoving frame the transfer equation has an extra term from the
frequency shift.  CMFGEN absorbs it into a \textbf{generalised opacity},
\begin{equation}
  \hat\chi \;=\; \chi \;+\; \frac{\Gamma}{\Delta\log\nu} ,
\end{equation}
where $\Delta\log\nu$ is the frequency-grid spacing.  In the standard
(monotonic) solver $\Gamma > 0$ so $\hat\chi > \chi > 0$.  In
non-monotonic flow $\Gamma$ can be negative enough to make
$\hat\chi < 0$, which is unphysical.  The Knop et al.\ (2009) fix
replaces $\Gamma$ with $|\Gamma|$ in the opacity (so $\hat\chi > 0$
always) and carries the sign in the source function instead.

\subsection{Summary table}

\begin{center}
\begin{tabular}{lll}
\toprule
Symbol & Name & Units \\
\midrule
$I_\nu$       & specific intensity        & $\mathrm{erg\,s^{-1}\,cm^{-2}\,Hz^{-1}\,sr^{-1}}$ \\
$J_\nu$       & mean intensity (0th moment) & same as $I_\nu$ \\
$H_\nu$       & first moment (flux)       & same as $I_\nu$ \\
$K_\nu$       & second moment (rad.\ pressure) & same as $I_\nu$ \\
$N_\nu$       & third moment              & same as $I_\nu$ \\
$\bar J$      & profile-averaged $J$      & same as $I_\nu$ \\
$\eta_\nu$    & emissivity                & $\mathrm{erg\,s^{-1}\,cm^{-3}\,Hz^{-1}\,sr^{-1}}$ \\
$\chi_\nu$    & opacity                   & $\mathrm{cm^{-1}}$ \\
$S_\nu$       & source function ($=\eta/\chi$) & same as $I_\nu$ \\
$\tau_\nu$    & optical depth             & dimensionless \\
$\hat\chi$    & generalised opacity       & $\mathrm{cm^{-1}}$ \\
$\hat\tau$    & generalised optical depth & dimensionless \\
$\mu$         & ray cosine                & dimensionless \\
$p$           & impact parameter          & cm \\
$\sigma$      & velocity-gradient parameter & dimensionless \\
$\Gamma$      & Sobolev frequency scale   & $\mathrm{s^{-1}}$ \\
$f_{\rm Edd}$ & Eddington factor ($K/J$)  & dimensionless \\
\bottomrule
\end{tabular}
\end{center}

\section{The two equations CMFGEN solves}
\label{sec:two-equations}

\subsection{1. Radiative transfer (the formal solution)}

At every point in the ejecta, the specific intensity $I_\nu$ obeys the
radiative transfer equation.  In the \textbf{comoving frame} (the frame
moving with the gas) this is, schematically,

\begin{equation}
  \frac{d I_\nu}{d s}
  \;=\; \eta_\nu \;-\; \chi_\nu\, I_\nu
  \;+\; \text{velocity-coupling terms},
\end{equation}

where $\eta_\nu$ is the emissivity, $\chi_\nu$ the opacity, and $s$ the
path length along a ray.  The ``velocity-coupling terms'' appear because
we are in an \emph{expanding} medium: in the comoving frame, the
frequency of a photon \emph{shifts} as it propagates.  In a monotonic
outflow photons are redshifted as they travel outward, coupling each
frequency to the next-lower frequency.  This coupling is the key
physical effect that distinguishes comoving-frame (CMF) radiative
transfer from static radiative transfer.

CMFGEN's standard formal solver absorbs the velocity coupling into a
\textbf{generalised opacity},

\begin{equation}
  \hat\chi \;=\; \chi \;+\; \frac{\Gamma}{\Delta\log\nu},
  \qquad
  \Gamma \;=\; \frac{V}{R}\bigl(1 + \sigma\,\mu^{2}\bigr),
  \qquad
  \sigma \;=\; \frac{d\ln V}{d\ln R} - 1,
\end{equation}

where $\Delta\log\nu$ is the frequency-grid spacing, $\mu$ the ray
cosine, and $\sigma$ the \emph{velocity-gradient parameter}.  The
transfer equation then becomes a standard short-characteristics problem,
\begin{equation}
  \frac{d I}{d\hat\tau} \;=\; I - S',
  \qquad
  \hat\tau = \hat\chi\, dz,
  \qquad
  S' = S + Q\, I_{\rm upwind},
  \qquad
  Q = \frac{\Gamma/\Delta\log\nu}{\hat\chi},
\end{equation}
with $S = \eta/\hat\chi$ the generalised source function and
$I_{\rm upwind}$ the intensity at the upwind frequency (already computed
in the decreasing-frequency sweep).  The formal solution is performed
for all rays (impact parameters $p$) and all frequencies, producing the
moments $J_\nu, H_\nu, K_\nu$ of the radiation field.

\subsection{2. Statistical equilibrium (rate equations)}

For every atom and every level, the population $n_i$ must satisfy
``number in = number out'':

\begin{equation}
  \sum_{j \neq i} n_j R_{ji} \;-\; n_i \sum_{j \neq i} R_{ij} \;=\; 0 .
\end{equation}

The rates $R_{ij}$ include:
\begin{itemize}[nosep]
  \item \textbf{Radiative rates}: $R_{ij}^{\rm rad} \propto \bar J_{ij}$,
        the mean intensity averaged over the line/continuum profile.
        These depend on the radiation field.
  \item \textbf{Collisional rates}: $R_{ij}^{\rm coll} \propto n_e$,
        depending on electron density and temperature.
\end{itemize}

So the populations depend on $\bar J$ (radiation), and $\bar J$ depends
on the populations (through $\eta_\nu, \chi_\nu$).  That is the coupling.

\section{The iteration scheme}

CMFGEN does not solve both equations simultaneously.  It alternates:

\begin{enumerate}[nosep]
  \item \textbf{Given current populations}, solve the radiative transfer
        equation to get the radiation field $J_\nu$ (and the Eddington
        factors, luminosity, etc.).
  \item \textbf{Given $J_\nu$}, solve the rate equations to get updated
        populations.
  \item Repeat until the populations stop changing.
\end{enumerate}

This is called \textbf{$\Lambda$-iteration} when step~1 uses a simple
formal solution.  Plain $\Lambda$-iteration converges very slowly, so
CMFGEN uses several accelerations:

\begin{description}[nosep]
  \item[NG acceleration] (Ng 1974): every few iterations, extrapolate
        the population sequence to try to jump ahead to the converged
        answer.  Works well when the last few iterations are converging
        smoothly.
  \item[BA iterations] (the ``Ba'' matrix, an approximate-Newton /
        Broyden-like correction): instead of just feeding $J_\nu$ back
        into the rate equations, linearise the \emph{coupled} transfer +
        rate system and solve for a correction to the populations that
        accounts for how the radiation field will respond.  This is much
        faster than pure $\Lambda$-iteration but more expensive per step
        and can be unstable.
\end{description}

A typical CMFGEN run cycles through blocks like
\[
  \text{LAMBDA} \;\to\; \text{LAMBDA} \;\to\; \text{NG}
  \;\to\; \text{FULL (BA)} \;\to\; \text{LAMBDA} \;\to\; \dots
\]
starting with cheap $\Lambda$ iterations to get close, then switching to
the more expensive BA/Full iterations to finish converging.  The
iteration strategy is controlled by the \texttt{IT\_SPECIFIER} block in
the input, which alternates between
\texttt{LAMBDA(FIXT)}, \texttt{FULL(FIXT)}, \texttt{FULL(FIXBA)}, and
\texttt{DO\_NG} commands.

\subsubsection{Current iteration strategy for SN2024pba}
\label{sec:iteration-strategy}

The SN2024pba production configuration uses the following iteration
strategy, controlled by the \texttt{single\_final} profile in
\texttt{prepare\_VADAT.py}:

\begin{enumerate}[nosep]
  \item \textbf{Cold start from LTE} (\texttt{LTE\_EST=T}): initial
        level populations are set to LTE values computed from the
        hydrodynamic temperature and electron density.  This provides a
        stable starting point that prevents the extreme population
        excursions that caused the earlier blow-up without
        \texttt{LTE\_EST} (see \S\ref{sec:caveats}).
  \item \textbf{Pure Lambda iterations} (\texttt{DO\_LAM\_IT=T}): the
        first iterations are pure Lambda iterations with
        \texttt{MAX\_LAM=$10^{5}$}, allowing large fractional population
        changes per iteration.  This is stable with the Knop formal
        solver because $\hat\chi > 0$ always, so $J_\nu$ remains finite
        even when the populations are far from converged.
  \item \textbf{Auto-switch to BA} (\texttt{DO\_LAM\_AUTO=T}): once the
        maximum population correction falls below 50\%, CMFGEN
        automatically switches from Lambda iterations to full BA
        (linearisation) iterations.  The BA solve couples the radiation
        field to the population correction, giving much faster
        convergence than pure Lambda iteration.
  \item \textbf{Ng acceleration} (\texttt{DO\_NG=T},
        \texttt{IBEG\_NG=50}, \texttt{BW\_NG=60}, \texttt{ITS/NG=8}):
        Ng acceleration is applied every 8 iterations once at least 50
        iterations have elapsed and the maximum correction is below
        \texttt{BEG\_NG}=10\%.  The Ng implementation includes
        per-depth rejection (if any population changes by more than
        $10.1\times$ or less than $0.09\times$), global cancellation
        (if more than 3 depths fail), and a post-Ng Lambda fallback.
        Ng is essential for the DIAG solver, which converges
        geometrically and would stall at $\sim 10^{10}$ without
        acceleration.
  \item \textbf{Trace-species freezing} (\texttt{FIX\_IMP=T}): trace
        species (low-abundance ions) are frozen in the BA matrix via
        \texttt{FIXPOP\_IN\_BA\_V3}, preventing singular BA matrix
        columns from trace species with near-zero populations.  This
        is particularly important for Fe\,\textsc{iv} and Fe\,\textsc{v}
        at late-time snapshots where their populations are set by
        photoionisation and can be extremely small.
  \item \textbf{Temperature handling} (\texttt{FIX\_T=T},
        \texttt{FIX\_T\_AUTO=T}): the orchestrator default fixes
        temperature to the hydrodynamic value at all depths, which is
        appropriate for a static snapshot.  The
        \texttt{FIX\_T\_AUTO=T} flag signals \texttt{DO\_LAM\_AUTO=T}
        (auto-switch from Lambda to BA).  Some configurations use
        \texttt{FIX\_T=F} with \texttt{FIX\_T\_AUTO=T} (Dessart et al.\
        2015 approach), where temperature is initially held fixed via
        the bootstrap while populations converge from LTE, then
        automatically released so CMFGEN solves for temperature
        self-consistently in the BA matrix.
\end{enumerate}

The convergence sequence is typically:
\[
  \text{LTE} \;\to\; \text{Lambda} \;\to\; \text{Lambda}
  \;\to\; \dots \;\to\; \text{BA (auto)} \;\to\; \text{Ng}
  \;\to\; \text{BA} \;\to\; \text{Ng} \;\to\; \dots
  \;\to\; \text{converged}
\]
with the Knop formal solver handling the compression zones at every
iteration so $J_\nu$ stays physical throughout.

\subsection{The moment equations and the BA matrix}

Instead of solving for $J_\nu$ directly from the ray quadrature, CMFGEN
\emph{optionally} solves the 0th and 1st moments of the transfer
equation as a tridiagonal system.  This couples adjacent depth points
and gives a more stable $J_\nu$ than pure ray quadrature, especially in
optically thick regions.  The moment equations require Eddington
factors ($K/J$ and $N/J$), computed from the ray moments.

The BA (Band Acceleration) method computes the Jacobian of $J_\nu$ with
respect to the populations ($dJ/d\chi$, $dJ/d\eta$) by linearising the
moment equations.  The resulting banded matrix couples adjacent depth
points; it is solved for a population correction that accounts for the
\emph{global} sensitivity of $J_\nu$ to population changes --- which is
why BA converges much faster than pure $\Lambda$ iteration.

\section{Why the velocity matters: the Sobolev / monotonic assumption}

In a standard supernova outflow, the velocity increases monotonically
outward: $V(r)$ increases with $r$, so $dV/dR > 0$ and
$\sigma = d\ln V/d\ln R - 1 > -1$.  This is the ``Hubble flow'' regime,
and it is hugely convenient because of the \textbf{Sobolev
approximation}: a photon at frequency $\nu$ only interacts with material
where the Doppler shift brings it into resonance, and if $dV/dR > 0$
that resonance point is \emph{unique}.  Each frequency maps to exactly
one depth.  The transfer problem becomes local and tractable.

Even when CMFGEN does a full formal solution (not the pure Sobolev
approximation), most of its machinery --- the characteristic rays, the
frequency-to-depth mapping, the way the source function is built up ---
\emph{assumes} that $V(r)$ is monotonic increasing outward.  The
characteristics don't cross, and the generalised opacity
$\hat\chi = \chi + \Gamma/\Delta\log\nu$ is always positive because
$\Gamma > 0$.

\section{The problem: non-monotonic velocity}

In interacting supernova systems --- SNe IIn, SNe with circumstellar
interaction, shock breakout, cold dense shells, reverse shocks --- the
velocity field can be \textbf{non-monotonic}:

\begin{itemize}[nosep]
  \item \textbf{Shocks}: $V$ jumps discontinuously at the shock front.
  \item \textbf{Velocity reversals}: $V$ decreases with $R$ in the
        post-shock region.
  \item \textbf{Cold dense shells}: ejecta decelerates abruptly, $V(R)$
        has a local minimum.
  \item \textbf{Reverse shocks}: inflowing material ($V < 0$) behind a
        forward shock.
\end{itemize}

In these cases $\sigma = d\ln V/d\ln R - 1$ can be $< -1$, meaning the
comoving-frame frequency gradient
$\Gamma = (V/R)(1 + \sigma\mu^{2})$ can be \emph{negative} for some ray
angles.

\subsection{Why the standard solver fails}

The standard CMFGEN formal solver was designed for monotonic outflows
where $\Gamma > 0$ always.  When $\Gamma < 0$ (non-monotonic flow), the
generalised opacity
\[
  \hat\chi \;=\; \chi \;+\; \frac{\Gamma}{\Delta\log\nu}
\]
can become \emph{negative}.  A negative opacity is unphysical and
causes a cascade of failures:

\begin{itemize}[nosep]
  \item \textbf{Negative optical depths} ($\hat\tau < 0$).
  \item \textbf{Exponentially growing intensities} ($e^{+\hat\tau}$
        instead of $e^{-\hat\tau}$).
  \item \textbf{NaN propagation} through the ray integration.
  \item \textbf{Instability} in the moment equations.
  \item \textbf{BA matrix singularity}.
\end{itemize}

\subsection{The infrastructure clamp (and why it is not enough)}

To prevent outright crashes, CMFGEN clamps $\sigma$ everywhere:
\[
  \texttt{IF}(\sigma < -0.999)\;\sigma \leftarrow -0.999 .
\]
This forces $\Gamma > 0$ (since $1 + (-0.999)\mu^{2} > 0$ for
$\mu \le 1$), so $\hat\chi > 0$ and the standard solver runs.  But the
clamp:

\begin{enumerate}[nosep]
  \item \textbf{Destroys the physics}: the true velocity gradient in
        compression zones is replaced by a flat $-0.999$ plateau.  The
        formal solution no longer represents the actual flow.
  \item \textbf{Suppresses resonance structure}: in non-monotonic flows,
        a single observer-frame frequency can resonate with
        \emph{multiple} comoving frequencies (multi-root resonance).
        The clamp eliminates this.
  \item \textbf{Corrupts spectra}: lines formed in compression zones
        (shock-cooled material, cold dense shells) get wrong profiles
        because the velocity coupling is wrong.
  \item \textbf{Breaks energy balance}: the wrong $J_\nu$ feeds into the
        rate equations, giving wrong populations and wrong temperatures.
\end{enumerate}

\subsection{The moment-equation and BA-matrix problems}

Even if the formal solver could handle $\sigma < -1$, the moment
equations contain terms like $(1 + \sigma)$ that become large and
negative when $\sigma < -1$.  The tridiagonal solver becomes unstable
because the diagonal terms can change sign, the off-diagonal coupling
terms reverse direction, and the Thomas algorithm assumes diagonal
dominance.  The same problem affects the BA matrix: the BA variation
uses $(1 + \sigma)$ terms, so with $\sigma < -1$ the BA matrix can
become singular or ill-conditioned.

\subsection{Summary of the problem}

The standard CMFGEN solver cannot handle non-monotonic velocity fields
because:
\begin{enumerate}[nosep]
  \item The generalised opacity
        $\hat\chi = \chi + \Gamma/\Delta\log\nu$ becomes negative when
        $\Gamma < 0$, causing numerical instability.
  \item The infrastructure clamp ($\sigma \ge -0.999$) prevents crashes
        but destroys the physics in compression zones.
  \item The moment equations and BA matrix are unstable for
        $\sigma < -1$.
  \item Multi-root resonance (multiple comoving frequencies coupling to
        one observer frequency) is not handled.
\end{enumerate}

\section{The non-monotonic solver (Knop, Hauschildt \& Baron 2009)}

\subsection{Reference}

The fix is a special formal solver that implements the formalism of

\begin{quote}
  Knop, S., Hauschildt, P. H., \& Baron, E.\ 2009,
  ``A non-monotonic velocity law for the radiative transfer in
  supernova atmospheres'', \emph{A\&A}, 493, 1197.
\end{quote}

In the codebase this is the routine
\texttt{FG\_J\_CMF\_NONMONO\_V1}.  It only changes \emph{how $J_\nu$ is
computed in the formal solution}; the rest of the CMFGEN iteration cycle
(rate equations, temperature correction, NG acceleration, Sobolev
escape probabilities, line opacities, atomic data input) is unchanged.

\subsection{A simple analogy: the reversing river}
\label{sec:analogy}

Before the mathematics, here is an analogy that captures the whole idea.

\medskip\noindent
\textbf{The setup.}\quad
Imagine a river flowing through a series of towns.  The river is the
\emph{frequency axis}: each town sits at one frequency.  Leaves floating
on the river are \emph{photons}.  The current is the
\emph{velocity coupling} $\Gamma$: it carries leaves from one town to
the next.  The muddiness of the water is the \emph{opacity} $\hat\chi$:
it determines how far a leaf can travel before it sinks (is absorbed) or
is replaced by a new one (is emitted).

\medskip\noindent
\textbf{Normal flow (monotonic expansion).}\quad
In a normal river the current always flows \emph{downstream} --- from
high frequency to low frequency (redshift).  Each town receives its
leaves from the town \emph{upstream} (the next-higher frequency, which
was already visited in the sweep).  The muddiness is always positive
(the current adds to it), so leaves travel a finite distance and
everything is well-behaved.  This is the standard CMFGEN solver.

\medskip\noindent
\textbf{The problem: a reversing river.}\quad
Now imagine the river hits a gorge and the current \emph{reverses} in
some stretches --- the water flows \emph{upstream} there.  This is a
compression zone ($\Gamma < 0$).  Two things go wrong:

\begin{enumerate}[nosep]
  \item \textbf{Negative muddiness.}  The old formula added the current
        to the muddiness.  When the current reverses, it
        \emph{subtracts}, and the muddiness can go negative.  A river
        with negative muddiness is nonsensical --- leaves would
        \emph{grow} stronger as they travel, exponentially.  In the
        code this means $e^{+\hat\tau}$ instead of $e^{-\hat\tau}$:
        intensities blow up or go negative, and NaNs propagate.
  \item \textbf{Wrong upstream town.}  In the reversed stretch, leaves
        now arrive from \emph{downstream} (the next-\emph{lower}
        frequency), not from upstream.  But the standard solver always
        asks the upstream town for its leaves --- so it asks the wrong
        neighbour and gets the wrong answer.
\end{enumerate}

\medskip\noindent
\textbf{The Knop fix, in river terms.}\quad
Knop et al.\ (2009) make two changes:

\begin{enumerate}[nosep]
  \item \textbf{Muddiness is always positive.}  Use the
        \emph{strength} of the current (its absolute value $|\Gamma|$)
        to set the muddiness, never its signed value.  The river is
        always muddy --- leaves always attenuate.  The muddiness no
        longer cares which way the current flows.

  \item \textbf{Ask the right neighbour.}  Keep the \emph{signed}
        current in a separate quantity (the source-function correction
        $Q$).  When $Q > 0$ (flowing downstream), ask the upstream town
        (the previous frequency, already computed).  When $Q < 0$
        (flowing upstream), ask the \emph{downstream} town (the next
        frequency).  The sign of $Q$ tells each town which neighbour to
        listen to.
\end{enumerate}

\medskip\noindent
\textbf{The catch: the downstream town hasn't spoken yet.}\quad
CMFGEN visits towns in \emph{downstream} order (high frequency to low
frequency).  In a reversed stretch, the town you need to ask
(downstream, lower frequency) hasn't been visited \emph{yet} in this
round.  So on the first pass you pretend it said nothing
($I_{\rm NEXT} = 0$) --- a crude but stable guess.  After each complete
round (a ``great iteration''), you \emph{remember} what every town said
and use that memory next time.  Over several rounds the remembered
answers converge to the truth.  This is the Gauss--Seidel scheme across
great iterations.

\medskip\noindent
\textbf{The dual-current bookkeeping.}\quad
To avoid confusing the rest of the CMFGEN machinery (which was built for
a one-way river), the solver keeps \emph{two} current maps:
\texttt{SIGMA\_TRUE} (the real, reversing current --- seen only by the
new solver) and \texttt{SIGMA} (a fake, always-downstream current
clamped to $\ge -0.999$, seen by everything else).  The old machinery
keeps working as if the river never reverses; only the new formal
solver knows the truth.

\medskip\noindent
\textbf{Summary of the analogy.}\quad

\begin{center}
\begin{tabular}{ll}
\toprule
River analogy & CMFGEN quantity \\
\midrule
River (towns along it)        & frequency grid \\
Leaves on the water           & photons ($I_\nu$) \\
Current                       & velocity coupling $\Gamma$ \\
Muddiness of the water        & generalised opacity $\hat\chi$ \\
Where leaves come from        & source function $S'$ \\
Upstream town                 & previous frequency ($I_{\rm PREV}$) \\
Downstream town               & next frequency ($I_{\rm NEXT}$) \\
Reversed stretch              & compression zone ($\Gamma < 0$) \\
One complete round of visits  & one great iteration \\
Remembering last round's answers & Gauss--Seidel coupling \\
Two current maps              & \texttt{SIGMA\_TRUE} vs \texttt{SIGMA} \\
\bottomrule
\end{tabular}
\end{center}

The key insight is that \emph{resistance} (opacity) and
\emph{direction} (which neighbour to ask) are two different things.  The
standard solver bundled them together; Knop et al.\ separated them.

\subsection{Core idea: separate opacity and source}

The key insight of Knop et al.\ (2009) is to use the
\textbf{magnitude} of $\Gamma$ in the opacity (keeping it positive)
while retaining the \textbf{signed} $\Gamma$ in the source function:

\begin{align}
  \hat\chi &= \chi \;+\; \frac{|\Gamma|}{\Delta\log\nu}
            &&\text{(always $> 0$)} \\
  Q        &= \frac{\Gamma/\Delta\log\nu}{\hat\chi}
            &&\text{(signed, can be negative)} \\
  S        &= \frac{\eta}{\hat\chi} \\
  S'       &= S \;+\; Q\, I_{\rm upwind}
            &&\text{(signed source)}
\end{align}

This ensures:
\begin{itemize}[nosep]
  \item $\hat\chi$ is always positive --- no negative optical depths.
  \item The signed $Q$ carries the \emph{physical direction} of
        frequency coupling.
  \item The source function $S'$ accounts for whether the flow is
        expanding ($Q > 0$) or compressing ($Q < 0$).
\end{itemize}

The transfer equation keeps the same form,
\[
  \frac{dI}{d\hat\tau} = I - S',
  \qquad \hat\tau = \hat\chi\, dz,
\]
which is always well-defined.  The physics of the compression is now
captured by the \emph{sign} of $Q$ and the choice of upwind intensity,
not by a negative opacity.

\subsection{Sign-dependent upwind frequency}

The \textbf{critical innovation} is the sign-dependent choice of upwind
frequency.

In the standard (monotonic) solver, the upwind frequency is always the
\emph{previous} frequency (the higher frequency, already computed in the
decreasing-frequency sweep).  This is correct for expanding flow where
$\Gamma > 0$: a photon at frequency $\nu$ couples to the next-higher
frequency $\nu + d\nu$ (redshifted as it travels outward).

In non-monotonic flow, compressing regions ($\Gamma < 0$) \emph{reverse}
the frequency coupling direction: a photon at frequency $\nu$ couples to
the \emph{next-lower} frequency $\nu - d\nu$ (blueshifted as it travels
into the compression).  The upwind frequency is therefore the
\emph{next} frequency, not the previous one.  The solver selects the
upwind intensity based on the sign of $\Gamma$:

\begin{equation}
  I_{\rm upwind} \;=\;
  \begin{cases}
    I_{\rm PREV} & \text{if } \Gamma \ge 0 \quad\text{(expanding: standard)} \\
    I_{\rm NEXT} & \text{if } \Gamma < 0 \quad\text{(compressing: Knop fix)}
  \end{cases}
\end{equation}

\subsection{Gauss--Seidel frequency coupling across great iterations}

There is a catch: CMFGEN sweeps frequencies in \emph{decreasing} order.
At frequency $\nu_i$, the next frequency $\nu_{i+1}$ has \emph{not yet}
been computed in the current sweep, so $I_{\rm NEXT}$ is not available
directly.

The Knop et al.\ solution is a \textbf{Gauss--Seidel approach across
great iterations}:

\begin{itemize}[nosep]
  \item On the \emph{first} great iteration, $I_{\rm NEXT} = 0$ (cold
        start).  Compressing points get $S' = S$ (no frequency
        coupling) --- a crude but stable approximation.
  \item After each great iteration, the computed intensities at each
        frequency and depth are \textbf{stored} for use as
        $I_{\rm NEXT}$ in the \emph{next} great iteration.
  \item Over successive great iterations, $I_{\rm NEXT}$ converges to
        the correct value, and the frequency coupling in compression
        zones becomes physically accurate.
\end{itemize}

This is directly analogous to how $\Lambda$ iteration converges: each
iteration uses the previous iteration's radiation field, refining it
until convergence.

\subsection{The dual-$\sigma$ architecture}

To preserve compatibility with the existing CMFGEN infrastructure, the
solver maintains \emph{two} $\sigma$ arrays:

\begin{center}
\begin{tabular}{lll}
\toprule
Array & Value & Used by \\
\midrule
\texttt{SIGMA\_TRUE} & unclamped $d\ln V/d\ln R - 1$ &
   the non-monotonic formal solver only \\
\texttt{SIGMA}       & clamped ($\ge -0.999$)        &
   all other routines: moment equations, \\
                     &                              &
   BA matrix, NG acceleration, Sobolev \\
                     &                              &
   escape probabilities, regridding, \\
                     &                              &
   grey opacity \\
\bottomrule
\end{tabular}
\end{center}

This separation ensures:
\begin{itemize}[nosep]
  \item The formal solution sees the \emph{true} non-monotonic physics.
  \item The infrastructure routines remain \emph{stable}.
  \item Existing models with monotonic velocity are completely
        unaffected --- if there are no compressing depths, the solver
        behaves identically to the standard monotonic path.
\end{itemize}

\texttt{SIGMA\_TRUE} is set wherever the velocity law is established
(the hydrodynamical-model setup, the SN IIn velocity-law routine, the
OBS module, and on warm restart), and is extended to the ray grid by
linear interpolation (not in log space, because $\log(\sigma+1)$ is
undefined for $\sigma < -1$).

\subsection{Bypassing the moment equations}

Because the moment equations are unstable for $\sigma < -1$, the
non-monotonic solver must \emph{bypass} them.  This is controlled by the
VADAT flag \texttt{USE\_LAM\_ES}:

\begin{itemize}[nosep]
  \item \texttt{USE\_LAM\_ES = F}: solve the moment equations
        (standard path, unstable for $\sigma < -1$).
  \item \texttt{USE\_LAM\_ES = T}: use the ray-based $J$ directly, and
        compute Eddington factors from the ray moments via
        \texttt{GET\_MOMS\_NON\_REL} (the non-monotonic path).
\end{itemize}

The Eddington factors ($K/J$, $N/J$) computed from the ray moments are
still physically correct --- they come from the true formal solution ---
they just bypass the unstable tridiagonal moment solve.

\subsection{The BA approximation in non-monotonic zones}

The BA variation (\texttt{DO\_VAR\_CONT} $\to$
\texttt{VAR\_MOM\_J\_CMF\_V12}) uses the \emph{clamped} $\sigma$
($\ge -0.999$), not \texttt{SIGMA\_TRUE}.  This means the BA Jacobian is
\emph{approximate} in compression zones --- it does not capture the full
non-monotonic physics.  This is a deliberate trade-off:

\begin{itemize}[nosep]
  \item The formal solution ($J$ computation) uses the true physics.
  \item The BA matrix (Jacobian) uses a clamped-$\sigma$ approximation.
  \item With \texttt{DIAG} (\texttt{NUM\_BNDS=1}), the BA iteration is
        stable but may converge more slowly in strongly non-monotonic
        regions.
  \item With \texttt{TRIDIAG} (\texttt{NUM\_BNDS=3}), the BA matrix can
        go \emph{singular} (see \S\ref{sec:caveats} and the
        ``Required configuration'' section below).
\end{itemize}

The Eddington factors used by the BA variation \emph{are} computed from
the true formal solution (via \texttt{GET\_MOMS\_NON\_REL}), so the BA
operator has the correct radiation-field geometry even though its
velocity coupling is approximate.  In practice, the iteration strategy
compensates by alternating LAMBDA iterations (stable, no BA),
\texttt{FULL(FIXT)} iterations (BA with approximate Jacobian), and NG
acceleration (extrapolation).

\subsection{Required configuration}
\label{sec:required-config}

The non-monotonic solver is activated by setting
\texttt{NON\_MONO\_SOLVER=T} in the VADAT file.  When this flag is
present, the \texttt{enforce\_sn\_iin\_physics()} routine in
\texttt{prepare\_VADAT.py} automatically forces a set of
\emph{non-negotiable} settings that are required for the Knop formal
solver to function correctly.  These forced settings cannot be
overridden by the user or by the profile defaults, because incorrect
values cause crashes or physically inconsistent results:

\begin{center}
\begin{tabular}{lll}
\toprule
Setting & Forced value & Reason \\
\midrule
\texttt{NON\_MONO\_SOLVER} & \texttt{T}  & Enable the Knop et al.\ (2009) formal solver \\
\texttt{GLOBAL\_LINE}      & \texttt{BLANK} & Blanketing mode (lines in continuum frequency set) \\
\texttt{DO\_SOB\_LINES}    & \texttt{F}  & BLANK handles lines; no separate Sobolev treatment \\
\texttt{LAM\_SET}          & \texttt{F}  & Do not switch to SOB for long wavelengths \\
\texttt{PURE\_HUB}         & \texttt{F}  & Preserve non-monotonic velocity from STELLA \\
\texttt{VEL\_LAW}          & \texttt{11} & Read SN hydrodynamic structure directly \\
\texttt{USE\_J\_REL}       & \texttt{F}  & No relativistic moment solver \\
\texttt{NEG\_OPAC\_OPT}    & \texttt{SRCE\_CHK} & Handle negative opacity via source check \\
\texttt{TRANS\_*}          & \texttt{BLANK} & All per-ion transitions through blanketing set \\
\bottomrule
\end{tabular}
\end{center}

Note that \texttt{GLOBAL\_LINE} is set to \texttt{SOB} by both the
\texttt{single\_final} profile in \texttt{prepare\_VADAT.py} and by
\texttt{run\_model\_grid.py}'s \texttt{enforce\_cmfgen\_sn\_iin\_settings()},
but the final \texttt{enforce\_sn\_iin\_physics()} call overwrites it to
\texttt{BLANK} when \texttt{NON\_MONO\_SOLVER=T}.  The effective value is
therefore \texttt{BLANK}, but the provenance is the final enforcer, not
the profile.

The inner boundary condition (\texttt{DIF}, \texttt{IB\_METH}) is
\emph{not} forced by the non-monotonic enforcement block.  The
orchestrator (\texttt{run\_model\_grid.py}) defaults to
\texttt{DIF=F} / \texttt{IB\_METH=ZERO\_FLUX} (optically thin core), but
\texttt{prepare\_VADAT.py} auto-detects the inner-boundary optical depth
from \texttt{SN\_HYDRO\_DATA} and switches to \texttt{DIF=T} /
\texttt{IB\_METH=DIFFUSION} when $\tau_{\rm inner} \ge 50$, where the
diffusion boundary condition is physically appropriate.  Both settings
are user-overridable via \texttt{vadat\_overrides} in the YAML config.

The following settings are \emph{not} forced by the non-monotonic
enforcement block, because they are iteration-strategy choices rather
than solver-consistency requirements.  The effective value of each
setting depends on a multi-stage pipeline:
\begin{enumerate}[nosep]
  \item \texttt{run\_model\_grid.py} sets orchestrator defaults via
        \texttt{enforce\_cmfgen\_sn\_iin\_settings()} (applied only if
        not already set in the YAML \texttt{vadat\_overrides}).
  \item \texttt{prepare\_VADAT.py} \texttt{apply\_profile()} sets
        \texttt{single\_final} profile defaults (may overwrite the
        orchestrator defaults for non-user-overridable keys).
  \item YAML \texttt{vadat\_overrides} are applied (user-specified values
        take precedence for user-overridable keys).
  \item \texttt{enforce\_sn\_iin\_physics()} forces the non-negotiable
        settings listed above.
\end{enumerate}
Because different model generations in the repository use different YAML
overrides deliberately, the table below lists the orchestrator default
(from \texttt{run\_model\_grid.py}) and notes where active configs differ:

\begin{center}
\begin{tabular}{llll}
\toprule
Setting & Orchestrator default & Reason & Config variation \\
\midrule
\texttt{USE\_LAM\_ES}  & \texttt{T}  & Bypass moment equations ($\sigma<-1$) & --- \\
\texttt{N\_TYPE}       & \texttt{N\_ON\_J} & Stable at velocity reversals & --- (forced) \\
\texttt{SOL\_METH}     & \texttt{DIAG} & Depth-independent BA (stable) & Some configs use TRIDIAG \\
\texttt{NUM\_BNDS}     & \texttt{1}  & Matches DIAG & 3 if TRIDIAG \\
\texttt{LTE\_EST}      & \texttt{T}  & LTE initial populations for cold start & --- \\
\texttt{FIX\_T}        & \texttt{T}  & T fixed to hydro (static snapshot) & Some configs use F \\
\texttt{FIX\_T\_AUTO}  & \texttt{T}  & Auto-release FIX\_T for BA & --- \\
\texttt{FIX\_NE}       & \texttt{T}  & Ne fixed to hydro (static snapshot) & Some configs use F \\
\texttt{FIX\_IMP}      & \texttt{T}  & Freeze trace species in BA & --- \\
\texttt{DO\_LAM\_IT}   & \texttt{T}  & Start in Lambda mode (stable cold-start) & --- \\
\texttt{DO\_LAM\_AUTO} & \texttt{T}  & Auto-switch Lambda$\to$BA at MAXCH $< 50\%$ & --- \\
\texttt{DO\_NG}        & \texttt{T}  & Ng acceleration (needed for DIAG) & Some configs use F \\
\texttt{IBEG\_NG}      & \texttt{50} & Min iterations before first Ng & Some configs use 15 \\
\texttt{BW\_NG}        & \texttt{60} & Ng bandwidth for DIAG & Some configs use 30 \\
\texttt{BEG\_NG}       & \texttt{0.1} & Start Ng when MAXCH $< 10\%$ & --- \\
\texttt{ITS/NG}        & \texttt{8}  & Iterations between Ng accelerations & --- \\
\texttt{MAX\_LAM}      & \texttt{$10^{5}$} & Max fractional pop.\ change / Lambda iter & Configs vary: 1.5--$10^{5}$ \\
\texttt{MAX\_LIN}      & \texttt{10} & Max linearisation steps per BA iteration & --- \\
\texttt{NUM\_LAM}      & \texttt{10} & Lambda iterations before BA check & Profile uses 2 \\
\texttt{BA\_CHK\_FAC}  & \texttt{0.05} & Hold BA fixed when MAXCH $< 5\%$ & --- \\
\texttt{VTURB}         & \texttt{20} km/s & Turbulent velocity (Dessart+2015) & --- \\
\texttt{COH\_ES}       & \texttt{T}  & Coherent ES for convergence & --- \\
\texttt{METHOD}        & \texttt{LOGMON} & Log-monotonic interpolation & --- \\
\texttt{INC\_SHOCK\_POWER} & \texttt{F} & Shock heating already in hydro $T(r)$ & --- \\
\texttt{THK\_CONT}     & \texttt{T}  & Thick continuum BC (Stockholm M15) & Some configs use F \\
\texttt{SCL\_SL\_OPAC} & \texttt{T}  & Scale SL line opacities (prevent spikes) & --- (forced) \\
\texttt{REV\_RGRID}    & \texttt{T}  & Runtime tau-based R-grid refinement & Most configs override to F \\
\bottomrule
\end{tabular}
\end{center}

\textbf{Why DIAG, not TRIDIAG.}\quad
The standard CMFGEN recommendation for monotonic flows is
\texttt{SOL\_METH=TRIDIAG} (\texttt{NUM\_BNDS=3}), which couples
adjacent depth points in the BA matrix and gives faster convergence.
However, \textbf{TRIDIAG is unstable with the non-monotonic solver for
strongly reversing flows}.

The root cause (documented in the BA iteration audit,
\texttt{BA\_AUDIT\_REPORT.txt}) is a mathematical inconsistency in the
BA variation operator:
\begin{itemize}[nosep]
  \item \texttt{VAR\_MOM\_J\_CMF\_V12} (the BA variation) uses the
        \emph{clamped} $\sigma \ge -0.999$ --- i.e.\ it assumes
        monotonic flow in its moment-equation assembly.
  \item But the Eddington factors ($K/J$, $N/J$) fed into the BA
        operator come from \texttt{GET\_MOMS\_NON\_REL}, which computes
        them from the \emph{non-monotonic} formal solve
        (\texttt{FG\_J\_CMF\_NONMONO\_V1}) --- i.e.\ they reflect the
        true reversing flow.
  \item The Eddington factors say ``the radiation field has this
        (non-monotonic) geometry''; the moment equations say ``the
        velocity coupling works this (monotonic) way''.  The resulting
        operator is inconsistent and can be \textbf{singular or
        ill-conditioned}, causing the tridiagonal solve
        (\texttt{CMF\_BLKBAND\_V3}) to return NaN.
  \item The NaN guard in \texttt{SOLVEBA\_V13} silently replaces NaN
        corrections with 0.0, so the BA iteration appears to ``converge''
        (zero changes) while actually having \emph{no effect}.
\end{itemize}

\texttt{DIAG} (\texttt{NUM\_BNDS=1}) avoids this by solving each depth
independently --- there is no cross-depth coupling matrix to go singular.
The trade-off is slower convergence: without the BA's global
$J_\nu$-to-population coupling, convergence is geometric and relies on
Ng acceleration to finish.  The production configuration
(\texttt{depth\_independent\_solver=True} in the YAML config, which is
the default) therefore uses \texttt{DIAG}.

\textbf{When TRIDIAG may work.}\quad
For \emph{mildly} non-monotonic flows (few compressing depths,
$\sigma$ only slightly below $-1$), the inconsistency is small and
TRIDIAG may remain well-conditioned.  The Stockholm reference models
(M15-He-O-Ca-evo) converge with TRIDIAG/\texttt{N\_ON\_J}/\texttt{NUM\_BNDS=3}
in $\sim$40 iterations.  But for strongly interacting SNe with deep
velocity reversals (the SN2024pba case), TRIDIAG produces a singular BA
matrix and must not be used.  The \texttt{depth\_independent\_solver}
flag in the YAML config controls this: \texttt{true} (default) $\to$
DIAG; \texttt{false} $\to$ TRIDIAG (experimental, for mild cases only).

\textbf{Temperature handling.}\quad
The orchestrator default is \texttt{FIX\_T=T} with
\texttt{FIX\_T\_AUTO=T}: temperature is fixed to the hydrodynamic
(STELLA) value at all depths, which is appropriate for a static
snapshot where the hydrodynamic temperature is the physically
consistent value.  The \texttt{FIX\_T\_AUTO=T} flag signals
\texttt{DO\_LAM\_AUTO=T} (auto-switch from Lambda to BA iterations).
Some model configurations (e.g.\ \texttt{X\_nomono}) use
\texttt{FIX\_T=F} with \texttt{FIX\_T\_AUTO=T}, following the Dessart
et al.\ (2015) approach where temperature is initially held fixed
(via the bootstrap) while populations converge from LTE, then
automatically released so CMFGEN solves for temperature
self-consistently in the BA matrix.  Both approaches are available;
the choice depends on whether the hydrodynamic temperature is trusted
as the final value (\texttt{FIX\_T=T}) or whether CMFGEN should refine
it (\texttt{FIX\_T=F}).

\textbf{Electron density and impurities.}\quad
The orchestrator default is \texttt{FIX\_NE=T}: the electron density is
fixed to the STELLA (Saha) value, which is appropriate for static
snapshot models where the hydrodynamic electron density is the
physically consistent value.  Some configurations (e.g.\
\texttt{X\_nomono\_v3\_vary\_ne}) use \texttt{FIX\_NE=F} to allow
CMFGEN to solve for the electron density self-consistently, which can
be important in the optically thin CSM where STELLA's Saha electron
density may be inaccurate.  However, BA iterations do not work with
free electron density, so \texttt{FIX\_NE=F} models are restricted to
pure Lambda iterations.  \texttt{FIX\_IMP=T} freezes trace species
(low-abundance ions) in the BA matrix via \texttt{FIXPOP\_IN\_BA\_V3},
preventing singular BA matrix columns that would otherwise arise from
trace species with near-zero populations.

\textbf{Ng acceleration.}\quad
Ng acceleration is \emph{enabled} (\texttt{DO\_NG=T}) in the current
configuration, with \texttt{IBEG\_NG=50} (at least 50 iterations before
the first Ng) and \texttt{BW\_NG=60} (bandwidth).  This is essential
because the DIAG solver converges geometrically --- without Ng, pure
Lambda iteration stalls at $\sim 10^{10}$ after 100+ iterations.  The
Ng implementation in \texttt{do\_ng\_accel\_band\_v4.f} includes robust
safety mechanisms: per-depth rejection if any population changes by more
than $10.1\times$ or less than $0.09\times$, global cancellation if more
than 3 depths fail, and a post-Ng Lambda fallback if corrections exceed
50\%.  Ng only fires when \texttt{MAXCH} $<$ \texttt{BEG\_NG} (10\%) and
\texttt{MAIN\_COUNTER} $\ge$ \texttt{IBEG\_NG} (50), by which point BA
mode is active and \texttt{FIXPOP\_IN\_BA\_V3} freezes trace species.

The DIFFERENCE method (\texttt{FG\_OPT=DIFF}) is rejected because it
does not construct $\hat\chi = \chi + |\Gamma|/\Delta\log\nu$ and is
unstable for $\sigma < -1$.  \texttt{N\_ON\_J} is required because at
velocity reversals the first moment $H$ passes through zero, making
$N/H$ (\texttt{G\_ONLY}) singular; $N/J$ is always well-defined since
$J > 0$.

\subsection{NaN safety guards}

Even with the correct integrals, numerical noise can produce slightly
negative intensities or (worse) NaNs, which on the next iteration would
propagate through the whole grid.  The solver therefore includes
comprehensive NaN guards at five points in the propagation chain:

\begin{enumerate}[nosep]
  \item \textbf{Ray intensity clamping} ($I_M, I_P \ge 10^{-30}$ if
        $< 0$ or NaN) --- catches NaN from source-function advection in
        compression zones.  Uses the IEEE-754 self-comparison trick
        ($x \neq x$ is true only for NaN) because Fortran's
        $x < 0$ test returns \texttt{FALSE} for NaN.
  \item \textbf{\texttt{JNU\_STORE} clamping} ($\ge 10^{-30}$ if
        $< 10^{-30}$ or NaN) --- prevents division by zero in
        $f_{\rm Edd} = K/J$ and NaN Eddington factors entering the BA
        matrix.
  \item \textbf{\texttt{I\_NEXT} storage guard} (NaN check before
        writing to the Gauss--Seidel arrays) --- prevents NaN
        contaminating the next great iteration.
  \item \textbf{TC/FEDD clamping in \texttt{COMP\_J\_BLANK}} (after the
        formal solve, before the moment solver) --- TC
        $\ge 10^{-30}$, FEDD $\ge 0$ with default $1/3$ if NaN.
  \item \textbf{$K/J$ division guard in \texttt{GET\_MOMS\_NON\_REL}} ---
        clamps \texttt{JNU\_STORE} and \texttt{KNU\_STORE} before
        computing $K/J$, preventing NaN $K/J$ from entering the BA
        matrix.
\end{enumerate}

All guards use OpenMP \texttt{ATOMIC} operations on the diagnostic
counters to avoid race conditions in the parallel ray loop.

\subsection{Sparse storage optimisation}

Storing intensities at \emph{all} depths, all rays, and all frequencies
for the Gauss--Seidel $I_{\rm NEXT}$ arrays would require enormous
memory: $N_D \times N_P \times N_{\rm CF} \times 8$ bytes per array,
which can reach tens of GB for realistic grids.

The optimisation exploits the fact that only depths where
$\sigma_{\rm true} < -1$ can ever have $\Gamma < 0$ (for any ray angle).
These are the ``compressing'' depths.  The solver:

\begin{enumerate}[nosep]
  \item Counts $N_{\rm comp}$ = number of depths with
        $\sigma_{\rm true} < -1$.
  \item Builds a sparse index \texttt{COMP\_DEPTH\_IDX} mapping the
        full-depth index to a sparse index $1 \ldots N_{\rm comp}$
        (or 0 if not compressing).
  \item Allocates \texttt{I\_P\_DEPTH\_NEXT} and
        \texttt{I\_M\_DEPTH\_NEXT} with dimensions
        $(N_{\rm comp}, N_P, N_{\rm CF})$ --- typically much smaller
        than $(N_D, N_P, N_{\rm CF})$.
\end{enumerate}

If there are no compressing depths
(\texttt{HAS\_COMPRESSING = .FALSE.}), the large arrays are not
allocated at all, and the solver behaves identically to the standard
monotonic path.  An additional in-place optimisation writes intensities
directly to the \texttt{NEXT} arrays during the frequency sweep,
eliminating the need for separate current-iteration arrays and the
initial copy/zero.

\subsection{Warm restart handling}

On warm restart (\texttt{NEWMOD = .FALSE.}), the routine that normally
sets \texttt{SIGMA\_TRUE} is not called, so \texttt{SIGMA\_TRUE} may be
zero from allocation.  The solver re-derives it from $V$ and $R$,
\[
  \sigma_{\rm true}(I) \;=\; \frac{R(I)}{V(I)}\,\frac{dV}{dR} \;-\; 1,
\]
using centred differences for interior points and one-sided differences
at the boundaries.  In the OBS module, \texttt{SIGMA\_TRUE} is always
re-derived from $V$ and $R$ because the restart file stores only the
clamped $\sigma$.

\subsection{Diagnostics}

On the first frequency of the first iteration the solver prints:
\begin{itemize}[nosep]
  \item the number of compressing depths ($N_{\rm comp}$);
  \item the number of depths with $\sigma < -1$;
  \item the $\sigma_{\rm true}$ range (min, max);
  \item the number of rays with multi-root resonances;
  \item the resonance turning-point count (min, max);
  \item confirmation of $\hat\chi$, signed $Q$, and sign-dependent
        upwind.
\end{itemize}
During the run it tracks (for an end-of-run summary) the count of
clamped negative/NaN intensities and the count of clamped NaN mean
intensities.

\section{Putting it together: what actually happens in a run}

\begin{enumerate}[nosep]
  \item A hydrodynamics code (e.g.\ STELLA) gives the structure ---
        density, velocity, temperature, composition vs radius --- at a
        chosen snapshot time.
  \item This is converted into \texttt{SN\_HYDRO\_DATA} (CMFGEN's input
        format, indexed outside-in: depth 1 = outermost shell).
  \item CMFGEN builds an atomic model (which ions, which levels) based
        on the physical conditions.
  \item It iterates: $\Lambda$ iterations to get close, then BA/Full
        iterations to converge.  At every iteration the non-monotonic
        solver handles the compression zones so $J_\nu$ stays physical.
  \item When the populations and luminosity stop changing, CMFGEN does a
        final formal solution and writes out the emergent spectrum via
        the OBS module.
\end{enumerate}

\section{The STELLA-to-CMFGEN conversion pipeline}
\label{sec:conversion}

The conversion from a STELLA snapshot to a CMFGEN \texttt{SN\_HYDRO\_DATA}
file is performed by \texttt{convert\_stella\_swd\_to\_cmfgen.py} and
is a critical step that determines the quality of the CMFGEN run.  The
current pipeline follows the Dessart et al.\ (2015) approach with
several project-specific adaptations.

\subsection{Tau-based truncation (Dessart+2015)}

The STELLA grid is first truncated to the region
$10^{-5} \le \tau_{\rm Rosseland} \le 30$, following Dessart et al.\
(2015).  The outer boundary is set by the NLTE optical depth
($\tau_{\rm min} = 0.01$), and the inner boundary by the
constant-kappa Rosseland optical depth with $\kappa = 0.34\;\mathrm{cm^2\,g^{-1}}$
(typical SN ejecta opacity).  The constant kappa is used for the inner
boundary because STELLA's NLTE kappa is too low in cold neutral gas
(giving $\tau_{\rm max} \sim 1$ instead of $\sim 240$ with
$\kappa = 0.34$), so $\tau_{\rm max} = 30$ would never be enforced with
the NLTE kappa.  The truncation is \emph{strict}: if too few points
remain, the grid is not relaxed --- the adaptive grid algorithm adds
points where gradients require them, but the $\tau = 30$ boundary is
preserved.

\subsection{Adaptive grid refinement (AMR)}

The truncated STELLA grid is remapped onto a CMFGEN depth grid using a
tau-based adaptive mesh refinement algorithm.  The algorithm:

\begin{enumerate}[nosep]
  \item Starts from the tau-trimmed STELLA grid as the base (typically
        80--150 points after truncation).
  \item Computes diagnostic gradients: radius, temperature, density,
        electron density, velocity, optical depth, and the velocity
        structure parameter $\sigma$.
  \item Identifies the worst-violating edge (highest score relative to
        its threshold) and inserts a midpoint.
  \item Repeats until all criteria are satisfied or \texttt{max\_points}
        is reached.
  \item Inserts \texttt{n\_outer\_extra}=5 additional points in the
        outermost low-tau region for better spectral accuracy.
  \item Runs a final grid+sigma smoothing pass to resolve any remaining
        violations.
\end{enumerate}

The current AMR criteria and thresholds are:

\begin{center}
\begin{tabular}{lll}
\toprule
Criterion & Threshold & Description \\
\midrule
\texttt{max\_radius\_jump}    & 0.05  & Max fractional radius jump $dr/r$ \\
\texttt{max\_temp\_jump}      & 0.10  & Max fractional temperature jump $dT/T$ \\
\texttt{max\_density\_jump}   & 0.15  & Max fractional density jump $d\rho/\rho$ \\
\texttt{max\_ne\_jump}        & 0.10  & Max fractional electron density jump \\
\texttt{max\_logtau\_jump}    & 0.10  & Max $\Delta\log\tau$ per cell \\
\texttt{max\_velocity\_jump}  & 0.10  & Max fractional velocity jump $dV/V$ \\
\texttt{max\_sigma\_jump}     & 1.0   & Max $|\Delta\sigma|$ per cell \\
\texttt{max\_cell\_ratio}     & 1.5   & Max ratio of adjacent cell sizes \\
\texttt{max\_points}          & 100--150 & Upper grid size limit \\
\texttt{min\_points}          & 60--80  & Minimum grid size (fill if below) \\
\bottomrule
\end{tabular}
\end{center}

Ionisation fronts are detected separately (via $|d\log n_e/d\log r| >
0.5$) and force-include a minimum of 5--10 source zones around each
front to ensure adequate resolution.  The AMR algorithm uses a
``worst-first'' insertion strategy: the single worst-violating edge is
refined at each step, guaranteeing that the highest-gradient regions
receive points first.

\subsection{Velocity smoothing}

Small numerical oscillations in the STELLA velocity profile
($|\Delta V|/V \sim 1\%$ over a few zones) are smoothed with a
Savitzky--Golay filter (window 5 zones, order 2) to prevent the
structure parameter $\sigma$ from reaching unphysical values
($\sigma \ll -100$) that would destabilise the CMFGEN Sobolev integrals.
The smoothing preserves the peaks of the velocity profile (unlike
Gaussian smoothing, which would reduce the shock velocity jump) by
redistributing radius rather than convolving values.  Velocity
smoothing can be disabled via the YAML config
(\texttt{smooth\_velocity: false}) for models where the STELLA velocity
is already smooth.

\subsection{Temperature profile}

The temperature profile from STELLA is used directly, with a floor of
$T_{\rm floor} = 2000$--$5000\;\mathrm{K}$ (configurable) to prevent
numerical issues in the cold outer CSM.  The Dessart et al.\ (2015)
adiabatic temperature profile is preserved in the conversion: the
STELLA gas temperature is used without modification, and the CSM
temperature is set to 6000\;K (typical photoionised CSM temperature)
where the STELLA profile drops below the floor.

\subsection{Composition and Ni56}

The STELLA 15-element composition is mapped onto CMFGEN's species list.
The $^{56}$Ni distribution is read from the STELLA \texttt{.xni} file
(Fortran unformatted \texttt{REAL*8}) and included in the
\texttt{SN\_HYDRO\_DATA} abundance table so that CMFGEN can compute
radioactive decay heating self-consistently.  For interaction-only
models with no $^{56}$Ni (the SN\_B case), a zero \texttt{.xni} file is
generated and the slurm template zeroes the \texttt{.xni} file after the
\texttt{eve2} pre-processor runs, ensuring no spurious radioactive
decay.

\section{Caveats and limitations of the non-monotonic solver}
\label{sec:caveats}

The Knop et al.\ (2009) solver makes the formal solution stable and
physically meaningful in compression zones, but it is \emph{not} a
complete non-monotonic radiative-transfer code.  This section lists the
known caveats so that users can judge when the results are reliable and
when a different tool (e.g.\ a ray-by-ray Lagrangian solver, or a code
such as SEDONA/SuperLite that does not assume monotonic flow) is more
appropriate.

\subsection{It is still a Sobolev / single-resonance approximation}

The solver operates within CMFGEN's comoving-frame formal-solution
framework, which inherits the Sobolev (large-velocity-gradient)
picture: each frequency is assumed to couple to \emph{one} upwind
frequency at each depth.  In a genuine velocity \emph{reversal} a single
observer-frame frequency can resonate with \emph{multiple} comoving
frequencies along a single ray (multi-root resonance).  The solver picks
a single upwind frequency based on the \emph{local} sign of $\Gamma$; it
does not trace all resonance points.  Lines formed deep inside a
strongly reversing flow can therefore have wrong profiles even though
$J_\nu$ stays finite.

\textbf{Implication:} reliable for shocks and decelerating shells where
the reversal is mild or where the line forms outside the reversal;
unreliable for lines that form \emph{inside} a deep velocity minimum
with multiple resonance points.

\subsection{The \texorpdfstring{$\sigma$}{sigma} clamp is numerical, not physical}

Every CMFGEN routine \emph{except} the non-monotonic formal solver sees
the clamped $\sigma \ge -0.999$.  This includes:
\begin{itemize}[nosep]
  \item the moment equations (\texttt{MOM\_J\_CMF\_V11},
        \texttt{MOM\_JREL\_V9});
  \item the BA / variation matrix (\texttt{VAR\_MOM\_J\_CMF\_V12},
        \texttt{VAR\_JREL\_V5});
  \item the Sobolev escape-probability routines (\texttt{SOBEW},
        \texttt{SOBJBAR}, \texttt{BETANEW});
  \item the grey opacity solve (\texttt{JGREY\_WITH\_FVT});
  \item NG acceleration and regridding.
\end{itemize}
The clamp creates a flat $\sigma = -0.999$ plateau wherever the true
$\sigma < -1$.  In those plateau regions the moment equations and the
BA Jacobian are \emph{approximate}: they use a well-conditioned but
unphysical velocity gradient.  The formal solution itself sees the true
$\sigma$ (via \texttt{SIGMA\_TRUE}), so $J_\nu$ is correct, but the
\emph{coupling} between $J_\nu$ and the population correction (the BA
step) is approximate in compression zones.

\textbf{Implication:} convergence in strongly non-monotonic regions may
be slower than in monotonic regions, because the BA operator is not the
true Jacobian there.  The iteration strategy (alternating LAMBDA,
FULL(FIXT), and NG) compensates, but a model with a large fraction of
compressing depths may need many more iterations.

\subsection{The BA matrix uses an approximate Jacobian in compression zones}

As noted above, the BA variation (\texttt{DO\_VAR\_CONT} $\to$
\texttt{VAR\_MOM\_J\_CMF\_V12}) uses the clamped $\sigma$, not
\texttt{SIGMA\_TRUE}.  This is a deliberate trade-off: the true
$\sigma < -1$ would make the BA matrix singular or ill-conditioned
(because the $(1 + \sigma)$ terms change sign), so the clamped value is
used to keep the linear solve well-posed.  The Eddington factors fed
into the BA operator \emph{are} computed from the true formal solution
(via \texttt{GET\_MOMS\_NON\_REL}), so the radiation-field geometry is
correct even though the velocity-coupling part of the Jacobian is
approximate.

\textbf{With DIAG (\texttt{NUM\_BNDS=1}, the production setting):} the
BA step is stable.  Each depth is solved independently, so there is no
cross-depth coupling matrix to go singular.  The BA correction is
approximate in compression zones (the Jacobian uses clamped $\sigma$),
so convergence there may be slower than in monotonic flow.  In
pathological cases the BA correction can be too small, requiring extra
LAMBDA or NG iterations to finish.

\textbf{With TRIDIAG (\texttt{NUM\_BNDS=3}, experimental):} the BA
matrix \emph{can go singular} for strongly non-monotonic flows.  The
cross-depth coupling in the tridiagonal system combines the clamped
$\sigma$ (monotonic moment equations) with non-monotonic Eddington
factors, producing an inconsistent operator.  In the SN2024pba tests
this caused \texttt{CMF\_BLKBAND\_V3} to return NaN, which the
\texttt{SOLVEBA\_V13} NaN guard silently replaced with 0.0 --- making
the BA iteration have zero effect while appearing to converge.  This is
why the production configuration uses DIAG.  TRIDIAG may work for mildly
non-monotonic flows (see the Stockholm M15 reference) but must not be
used for deep velocity reversals.

\subsection{Gauss--Seidel coupling converges slowly on the first great iterations}

The sign-dependent upwind frequency in compression zones requires
$I_{\rm NEXT}$, which is not yet computed in the current
decreasing-frequency sweep.  The solver uses a Gauss--Seidel scheme
\emph{across great iterations}: on the first great iteration
$I_{\rm NEXT} = 0$, so compression zones get $S' = S$ (no frequency
coupling).  The coupling is filled in only over successive great
iterations.

\textbf{Implication:} the \emph{first} great iteration of a
non-monotonic model is a cruder approximation than a monotonic model's
first iteration.  If the model is restarted from a cold start (LTE
populations) with strong compression, the early iterations may show
large population swings before the Gauss--Seidel coupling converges.
A warm restart from a monotonic model of the same snapshot reduces this
problem.

\subsection{Velocity plateaus and discontinuities are not handled well}

The Sobolev approximation assumes the velocity scale height is much
larger than the photon mean free path.  Two regimes violate this:
\begin{itemize}[nosep]
  \item \textbf{Velocity plateaus} ($dV/dR \approx 0$): $\sigma \to -1$,
        $\Gamma \to 0$, and the generalised opacity reduces to the
        static opacity.  The frequency coupling vanishes, so the
        comoving-frame advantage is lost and the line profile is
        unresolved.
  \item \textbf{Sharp discontinuities} (true shocks with
        $|\Delta V|/V \sim 1$ over one zone): the centred-difference
        $\sigma$ is dominated by interpolation artefacts, and the
        single-resonance assumption fails because the resonance point
        jumps discontinuously.
\end{itemize}
CMFGEN's grid should resolve shocks over several zones; the
\texttt{max\_sigma\_jump} AMR criterion in the conversion pipeline is
designed for this.  A shock that fits in a single grid cell will
produce unreliable line profiles regardless of the solver.

\subsection{Flux conservation has not been quantitatively verified}

The solver preserves $J_\nu$ positivity and finiteness, and the
Eddington factors come from the true formal solution, but a
\emph{quantitative} flux-conservation test (comparing the integrated
emergent luminosity to the input luminosity + radioactive + interaction
deposition) has not been performed for the non-monotonic path.  The
algebraic analysis shows that the Knop et al.\ formulation is
conservative, but numerical round-off, the $\sigma$ clamp, and the
Gauss--Seidel lag can introduce small flux errors in compression zones.

\textbf{Implication:} treat the absolute luminosity of lines formed in
compression zones as uncertain at the $\sim$few-percent level.  The
\emph{shape} of the spectrum and the relative strengths of lines formed
outside compression zones are reliable.

\subsection{Restart reproducibility is not guaranteed}

On warm restart (\texttt{NEWMOD = .FALSE.}), \texttt{SIGMA\_TRUE} is
re-derived from $V$ and $R$ by finite differences, because the restart
file stores only the clamped $\sigma$.  The re-derivation uses centred
differences on the \emph{current} grid, which may differ slightly from
the grid on which the original $\sigma$ was computed (if regridding
occurred).  This means a restarted run may not be bit-for-bit
reproducible with respect to the original run in the compression zones.

\subsection{Very small grids (\texorpdfstring{$N_D < 100$}{ND < 100})}

The non-monotonic solver inserts boundary points and requires enough
interior points to resolve the compression region with centred
differences.  On very small grids ($N_D < 100$) the boundary insertion
can leave too few interior points, the $\sigma$ finite differences
become noisy, and the formal solution can oscillate.  The AMR pipeline
targets $N_D \sim 150$--$300$ for interacting-SN models, which is safe.

\subsection{Extreme \texorpdfstring{$\sigma \ll -1$}{sigma << -1}}

When $\sigma$ is very negative (e.g.\ $\sigma = -7.5$ at a strong
reverse shock), the clamp plateau is wide and the discrepancy between
the true physics and the clamped BA Jacobian is large.  The formal
solution is still correct (it uses \texttt{SIGMA\_TRUE}), but the BA
iteration may stall or converge to a slightly wrong population
distribution in those zones.  For such extreme cases, a Monte Carlo
code (SEDONA, SuperLite) that does not rely on the Sobolev approximation
is a better cross-check.

\subsection{What is \emph{not} changed by the solver}

It is worth restating what the non-monotonic solver does \emph{not}
touch, so that users know where to look for unrelated problems:
\begin{itemize}[nosep]
  \item the rate equations (\texttt{SOLVE\_FOR\_POPS});
  \item the atomic data input and level solving;
  \item the temperature correction (\texttt{DIFFW} / radiative
        equilibrium);
  \item NG acceleration (uses clamped $\sigma$, unchanged);
  \item Sobolev escape probabilities (Dessart+2015 guard, separate from
        the Knop solver);
  \item line opacities and continuum opacities;
  \item the OBS module (spectrum post-processing).
\end{itemize}
A convergence problem or a wrong line ratio that persists with the
non-monotonic solver off is \emph{not} caused by the non-monotonic
solver and should be diagnosed with the standard CMFGEN toolkit.

\subsection{When to use a different code}

The non-monotonic solver is the right tool when:
\begin{itemize}[nosep]
  \item the velocity field is mostly monotonic with a mild, resolved
        compression zone (shock, decelerating shell);
  \item the scientifically important lines form \emph{outside} the
        compression zone;
  \item a detailed non-LTE spectrum with many explicit levels is needed
        (CMFGEN's strength).
\end{itemize}
A different tool is preferable when:
\begin{itemize}[nosep]
  \item the flow has a deep, multi-resonance velocity reversal and the
        lines of interest form inside it --- use SEDONA or SuperLite
        (Monte Carlo / Lagrangian, no Sobolev assumption);
  \item the shock is unresolved ($< 3$ zones) --- refine the grid or use
        a code with shock-capturing hydrodynamics coupled to the
        transport;
  \item time-dependent radiation transport matters (e.g.\ shock
        breakout) --- CMFGEN is a snapshot code; use STELLA or a
        time-dependent RHD code.
\end{itemize}

\section{Key takeaways}

\begin{itemize}[nosep]
  \item CMFGEN solves \emph{coupled} radiative transfer + statistical
        equilibrium by \emph{iteration}, accelerated with NG and BA
        methods.
  \item Almost all of its machinery assumes the velocity increases
        monotonically outward (Sobolev picture), so the generalised
        opacity $\hat\chi = \chi + \Gamma/\Delta\log\nu$ is always
        positive.
  \item Interacting supernovae have \emph{compression zones} where
        $dV/dR < 0$, breaking that assumption and making
        $\sigma < -1$, so $\Gamma$ can be negative and $\hat\chi$ can go
        negative.
  \item The infrastructure clamp ($\sigma \ge -0.999$) prevents crashes
        but destroys the physics in compression zones.
  \item The Knop et al.\ (2009) non-monotonic solver fixes this by:
    \begin{enumerate}[nosep]
      \item using $|\Gamma|$ in the opacity (so $\hat\chi > 0$ always);
      \item retaining the signed $\Gamma$ in the source function (so the
            physical direction of frequency coupling is preserved);
      \item choosing the upwind frequency from the sign of $\Gamma$
            (previous frequency for expanding flow, next frequency for
            compressing flow);
      \item using a Gauss--Seidel scheme across great iterations to
            supply the not-yet-computed $I_{\rm NEXT}$ in compression
            zones;
      \item maintaining a dual-$\sigma$ architecture so the formal
            solver sees the true physics while the infrastructure
            routines stay stable;
      \item bypassing the unstable moment equations
            (\texttt{USE\_LAM\_ES = T}) and computing Eddington factors
            directly from the ray moments;
      \item using the clamped $\sigma$ in the BA matrix with
            \texttt{DIAG} (\texttt{NUM\_BNDS=1}) for stability --- the
            depth-independent solve avoids the singular cross-depth
            coupling that TRIDIAG produces when non-monotonic Eddington
            factors meet clamped-$\sigma$ moment equations;
      \item guarding against NaNs and negative intensities at five
            points in the propagation chain.
    \end{enumerate}
  \item It only needs to run on the compression depths, so a sparse
        storage optimisation keeps memory manageable.
  \item The rest of CMFGEN (rate equations, temperature correction, NG
        acceleration, Sobolev probabilities, line opacities, atomic
        data) is unchanged --- the solver is a drop-in replacement for
        the formal-solution step.
  \item The current production configuration (SN2024pba) uses:
    \begin{itemize}[nosep]
      \item LTE cold start (\texttt{LTE\_EST=T}) for stable initial
            populations;
      \item pure Lambda iterations with auto-switch to BA
            (\texttt{DO\_LAM\_IT=T}, \texttt{DO\_LAM\_AUTO=T});
      \item Ng acceleration enabled (\texttt{DO\_NG=T},
            \texttt{IBEG\_NG=50}, \texttt{BW\_NG=60}) --- essential for
            DIAG convergence;
      \item temperature fixed to hydro with auto-release
            (\texttt{FIX\_T=T}, \texttt{FIX\_T\_AUTO=T}; some configs
            use \texttt{FIX\_T=F} for self-consistent T, Dessart+2015);
      \item electron density fixed to hydro (\texttt{FIX\_NE=T};
            some configs use \texttt{FIX\_NE=F} for self-consistent
            NLTE ionisation, but this restricts to pure Lambda);
      \item trace species frozen in BA (\texttt{FIX\_IMP=T}) to prevent
            singular matrix columns;
      \item turbulent velocity 20\;km/s (\texttt{VTURB=20},
            Dessart+2015);
      \item coherent electron scattering (\texttt{COH\_ES=T}) for
            convergence;
      \item blanketing line treatment (\texttt{GLOBAL\_LINE=BLANK},
            forced by the non-monotonic enforcement when
            \texttt{NON\_MONO\_SOLVER=T});
      \item tau-based AMR grid with $\tau_{\rm Rosseland} \le 30$
            truncation (Dessart+2015) and ionisation-front-aware
            refinement;
      \item inner boundary auto-detected: \texttt{ZERO\_FLUX}
            (\texttt{DIF=F}) for $\tau_{\rm inner} < 50$, or
            \texttt{DIFFUSION} (\texttt{DIF=T}) for
            $\tau_{\rm inner} \ge 50$.
    \end{itemize}
\end{itemize}

\end{document}
